\documentclass[11.6pt]{article}
\usepackage{graphicx}
\usepackage{amsmath,amssymb,amsthm}
\usepackage[style=apa, backend=biber]{biblatex}
\usepackage{indentfirst}
\addbibresource{reference.bib}
\usepackage{geometry}
\usepackage{float}



\title{\vspace{-4cm}\textbf{Prehistoric and Ancient Mathematics}}

\author{
 \textbf{Alexandra De Guzman} \\
   BS Economics \\
  Department of Economics, UPLB \\
  \texttt{aadeguzman10@up.edu.ph}
  \and
  \textbf{Dane Ambing} \\
  BS Applied Physics \\
  Institute of Physics \\
  \texttt{dpambing@up.edu.ph}
  \and
   \textbf{Iya Umali} \\
   BS Mathematics Education \\
   Institute of Mathematics, UPLB \\
  \texttt{irumali@up.edu.ph}
  \and
   \textbf{JC Esparto} \\
   BS Mechanical Engineering \\
   Department of Mechanical Engineering, UPLB \\
  \texttt{jecruz10@up.edu.ph} 
  \and
  \textbf{Kyle Martin Villagen} \\
   BS Computer Science \\
   Institute of Computer Science, UPLB \\
  \texttt{kavillage@up.edu.ph}
  \and
  \textbf{Mikhail Roa} \\
   BS Forestry \\
   College of Forestry and Natural Resources, UPLB \\
  \texttt{mroa1@up.edu.ph}
  \and
  \textbf{Simone Nisolada} \\
  BS Food Science and Technology \\
  Institute of Food Science and Technology, UPLB \\
  \texttt{sanisolada@up.edu.ph}
  \and
  \textbf{Zion Joseph S. Busto} \\
  BS Statistics \\
  Institute of Statistics, UPLB \\
  \texttt{zsbusto@up.edu.ph}
}



\date{\today}

\begin{document}

\maketitle
\pagenumbering{roman}
\newpage
\section*{Abstract}
This micro-research explores the development of mathematics in different eras, focusing on their numeral system, mathematical records, and practical use of mathematics. It highlights how mathematics developed in response to everyday needs such as agriculture, land measurement, and construction. The study also shows how different eras of mathematical practices contributed to later mathematical advancements.



\newpage


\tableofcontents
\listoffigures
\newpage
\pagenumbering{arabic}

\section{Chronology of Ancient Mathematics}

\subsection{Sumerian}

\subsubsection{Brief Histoy}
Sumer (a region of Mesopotamia, modern-day Iraq) was the birthplace of writing, the wheel,
agriculture, the arch, the plow, irrigation and many other innovations, and is often referred to
as the Cradle of Civilization.

Sumerian mathematics initially developed largely as a response to bureaucratic needs
when their civilization settled and developed agriculture (possibly as early as the 6th
millennium BCE) for the measurement of plots of land, the taxation of individuals, etc.

\subsubsection{Numeral System}
The Sumerians developed a sophisticated numeral system based on the sexagesimal
system, which is a base-60 system. This means that they used a base of 60 rather than the
base-10 system we use today. The sexagesimal system had practical advantages for the
Sumerians in calculations involving fractions and divisions. This system influenced later
mathematical traditions in the region.

\subsubsection{Cuneiform Tablets}
The Sumerians used cuneiform writing on clay tablets to record mathematical information.
These tablets contain a wealth of mathematical knowledge, including numerical
calculations, geometric principles, and problem-solving methods. The survival of these
tablets provides modern scholars with insights into ancient Sumerian mathematical
practices.

\subsubsection{Geometry}
Sumerians were among the earliest cultures to develop practical geometry. They used
geometric principles in land surveying, construction, and irrigation projects. The concept of
area measurement, especially in relation to agricultural land, was an essential aspect of
Sumerian geometry.


\newpage

\subsection{Egyptian}

\begin{itemize}
    \item Base-10 System: The Egyptians used a base-10 number system written in hieroglyphic numerals. Each symbol represents a unit following: 1, 10, 100, 1000, 10000, 100000, 1000000. 
    \item These base-10 hieroglyphs were seen in stone carvings in Karnak, dating back to 1500 BCE.
    \item  Fractions in ancient Egypt were limited to unit fractions in the form of 1/n. In numerical hieroglyphs, the unit fraction is represented as an “eye” symbol.
    \item  After the number system, it evolved into a much more compact form called the hieratic numerals. Similar to base-10, it has symbols assigned to a numerical value. 
 
\end{itemize}

\begin{figure}[H]
    \centering
    \includegraphics[width=0.75\linewidth]{egypt.JPG}
    \caption{Egyptian Numeral Hieroglyphs}
    \label{fig:placeholder}
\end{figure}


\begin{figure}[H]
    \centering
    \includegraphics[width=0.5\linewidth]{egyptiann.JPG}
    \caption{Hieratic Numerals}
    \label{fig:placeholder}
\end{figure}


\newpage

\subsection{Babylonian}
\subsubsection{Babylonian numerals}
\begin{figure}[H]
    \centering
    \includegraphics[width=0.75\linewidth]{baby.JPG}
    \caption{Babylonian numerals
}
    \label{fig:placeholder}
\end{figure}
\begin{itemize}
    \item A sexagecimal (base-60) numerical system. It is a more advanced form of the sexagecimal systems used by the previous Sumerian and Akkadian empires as it uses a place-value system.
    \item Written in clay tablets using cuneiform symbols.
    \item Influenced modern timekeeping (seconds in a minute, minutes in an hour) and measurements of angles (360 degrees in a circle).

\end{itemize}
\subsubsection{Arithmetic, Algebra, and Geometry
}
\begin{itemize}
    \item Babylonians wrote clay tablets composed of multiplication and division tables to aid in their arithmetic.
    \item They also wrote tablets for squares, cubes, square and cube roots, and sum of squares and cubes.
   \item Babylonians were known to solve algebraic equations. They implemented step-by-step procedures that were specific to certain problems (like algorithms) where they made extensive use of the aforementioned tablets. This was used in the calculation and measurement of land area.
    \item Babylonian geometry focused on practicality where they primarily measured lengths, areas, and volumes. They were able to compute for the area of right and isosceles triangles and isosceles trapezoids, as well as cylinders and some prisms.
    \item They are also credited for the idea of dividing a circle into 360 degrees.

\end{itemize}
\subsubsection{Plimpton 322 and the Pythagorean Theorem
}
\begin{figure}[H]
    \centering
    \includegraphics[width=0.75\linewidth]{pythagorean.JPG}
    \caption{Plimpton 322 and the Pythagorean Theorem}
    \label{fig:placeholder}
\end{figure}
\begin{itemize}
    \item The Plimpton 322 is a clay tablet that contains a mathematical table composed of rows of pythagorean triples.
    \item The tablet is dated to be from 1800 BCE which predates the major discoveries in geometry made by the Greeks.
    \item The tablet is believed to be evidence which suggests that Babylonians had knowledge about a precursor to the Pythagorean theorem and may have practiced trigonometry.

\end{itemize}


\newpage

\subsection{Ancient Indian}
In this era, mathematics was developed primarily to measure their infrastructures efficiently and to design drainage systems. Furthermore, their mathematics was also developed to create fire altars with precise shapes and structure.
\subsubsection{Indus Valley (Harappan) Civilization — c. 2600–1900 BCE}

\begin{itemize}
    \item Developed standardized weights and measurements and practical geometry.
    \item Mathematics in this era was used for city planning, construction, trade, and manufacturing
    \item Their infrastructures and drainage systems show that they were knowledgeable of measurements and proportions
    \item The history of Indian mathematics is hard to decipher since the Indus sceipt is hard to understand
    \item There is no way to determine their written mathematical system
\end{itemize}

\subsubsection{Vedict Period ( roughly 1500–500 BCE)}

\begin{itemize}
    \item Mathematical knowledge developed further, especially in geometry and measurement
    \item The Śulbasūtras provided rules for constructing squares, rectangles, circles, and religious fire altars.
    \item They were able to find ways to approximate the value of $\sqrt{2}$
    \item This mathematics developed mainly because Vedic religious rituals required precisely shaped and measured fire altars.
\end{itemize}



\newpage
\subsection{Ancient Chinese}

Empire during this period: Shang Dynasty (c. 1600–1046 BCE) and Zhou Dynasty (c. 1046–256 BCE).

The forms of mathematics they developed were practical (1) methods of counting and calculation, including the use of counting rods to represent numbers and perform arithmetic. They also developed (2) methods for working with positive and negative numbers, as well as the (3) Lo Shu Magic Square, a 3x3 square in which each row, column, and diagonal adds up to 15.

\begin{figure}[H]
    \centering
    \includegraphics[width=0.5\linewidth]{counting.JPG}
    \caption{Methods of Counting}
    \label{fig:placeholder}
\end{figure}

\begin{figure}[H]
    \centering
    \includegraphics[width=0.5\linewidth]{tangina.JPG}
    \caption{Methods for Working with Positive and Negative }
    \label{fig:placeholder}
\end{figure}

\begin{figure}[H]
    \centering
    \includegraphics[width=0.5\linewidth]{hoo.JPG}
    \caption{Lo Shu Magic Square }
    \label{fig:placeholder}
\end{figure}

Magic Squares are numerical arrangements where the sum of each row, column, and diagonal is the same. In various cultures, they believed to have mystical properties and bring good luck.



\newpage
\subsection{Ancient Greek}
“The laws of nature are but the mathematical thoughts of God.” - Euclid

\subsubsection{Proof and Logic, Axioms, and Postulates
}
\textbf{The Five Postulates}
\begin{itemize}
    \item \textbf{Postulate 1: }  A straight line segment can be drawn joining any two points.

    \item \textbf{Postulate 2: }Any straight line segment can be extended indefinitely.

    \item \textbf{Postulate 3:}  A circle can be drawn with any center and radius.

    \item \textbf{Postulate 4: }   All right angles are congruent.

    \item \textbf{Postulate 5 (The Parallel Postulate):} If a line intersecting two lines makes interior angles on one side summing to less than two right angles, those lines intersect on that side when extended.

\subsubsection{Geometric Proofs
}
\begin{itemize}
    \item \textbf{Pythagoras’ Theorem: 
} the square of the longest side (c, the hypotenuse) equals the sum of the squares of the other two sides (a and b).

\item \textbf{Thales’ Intercept Theorem
: }  if a triangle is drawn within a circle with the long side as a diameter of the circle, then the opposite angle will always be a right angle.

\subsubsection{Zeno's Paradox}
a half plus a quarter plus an eighth plus a sixteenth, etc, etc, to infinity will never quite equal a whole. 




\end{itemize}



    

\end{itemize}

\newpage
\subsection{Roman}


\newpage
\subsection{Mayan}

\subsubsection{Mayan Mathematics}
Mayan mathematics was an advanced number system developed for astronomy and calendar calculations in the Mayan society. The Mayan civilisation had settled in Central America from about 2000 BCE, while the Classic Period of Maya lasted from about 250 CE to 900 CE. \\

\noindent
Mayan numerals consisted of 3 symbols:
\begin{itemize}
    \item a shell shape, counting as 0; 
    \item a dot, counting as 1, and used as an indicator of place values; and
    \item a bar, counting as 5.

\end{itemize}

Mayan mathematics employed a vigesimal number system based on base 20, where vertically indicated place values were written based on powers of 20. This number system was used for numbers exceeding 19.
*In calendar calculations, the third place value counted as 360 instead of 400. 

\begin{figure}[H]
    \centering
    \includegraphics[width=0.75\linewidth]{mayan.JPG}
    \caption{Mayan Numerals}
    \label{fig:placeholder}
\end{figure}






\newpage
\nocite{indian}
\nocite{ancient_chinese}
\nocite{mayan}
\nocite{baby}
\nocite{egypt}
\nocite{martzloff2006history}
\nocite{reference_from_maam}
\setlength{\bibitemsep}{12pt}   
\printbibliography


\end{document}
